
\begin{remark}
    Last time we proved that if \(B\) and \(B'\) are two bases of a root system
    \(\Delta\), then there exists \(w\in W_B\) such that
    \[
        w(B')=B.
    \]
    We now prove that such an element is unique.
\end{remark}

\begin{proposition}
    Let \(w\in W_B\). Then
    \[
        w(B)=B
        \qquad\Longleftrightarrow\qquad
        w=1.
    \]
\end{proposition}

\begin{proof}
    The implication \(w=1\Rightarrow w(B)=B\) is clear. Conversely, suppose
    that \(w(B)=B\). We prove that \(w=1\).

    Assume, to the contrary, that \(w\neq 1\). Write
    \[
        w=S_{\beta_1}S_{\beta_2}\cdots S_{\beta_m},
        \qquad \beta_i\in B,
    \]
    as a shortest expression. Thus \(m\geq 1\).

    We claim that
    \[
        w(\beta_m)\in \Delta^-.
    \]
    If \(m=1\), this is immediate, since
    \[
        S_{\beta_1}(\beta_1)=-\beta_1\in \Delta^-.
    \]
    Suppose now that \(m>1\), and suppose for contradiction that
    \[
        w(\beta_m)\in \Delta^+.
    \]
    Since
    \[
        S_{\beta_m}(\beta_m)=-\beta_m\in \Delta^-,
    \]
    we have
    \[
        S_{\beta_1}S_{\beta_2}\cdots S_{\beta_{m-1}}(\beta_m)\in \Delta^-.
    \]
    Consider the sequence
    \[
        \beta_m,\quad
        S_{\beta_{m-1}}(\beta_m),\quad
        S_{\beta_{m-2}}S_{\beta_{m-1}}(\beta_m),\quad
        \ldots,\quad
        S_{\beta_1}S_{\beta_2}\cdots S_{\beta_{m-1}}(\beta_m).
    \]
    Its first term lies in \(\Delta^+\), while its last term lies in \(\Delta^-\).
    Hence there exists \(k\) such that
    \[
        S_{\beta_{k+1}}\cdots S_{\beta_{m-1}}(\beta_m)\in \Delta^+
    \]
    but
    \[
        S_{\beta_k}S_{\beta_{k+1}}\cdots S_{\beta_{m-1}}(\beta_m)\in \Delta^-.
    \]
    By the standard property of simple reflections, if
    \(\alpha\in \Delta^+\), \(\beta\in B\), and
    \(S_\beta(\alpha)\in \Delta^-\), then \(\alpha=\beta\). Therefore
    \[
        S_{\beta_{k+1}}\cdots S_{\beta_{m-1}}(\beta_m)=\beta_k.
    \]
    Set
    \[
        u=S_{\beta_{k+1}}\cdots S_{\beta_{m-1}}.
    \]
    Then
    \[
        \beta_k=u(\beta_m),
    \]
    and hence
    \[
        S_{\beta_k}=uS_{\beta_m}u^{-1}.
    \]
    It follows that
    \[
        \begin{aligned}
            w
             & =
            S_{\beta_1}\cdots S_{\beta_{k-1}}
            S_{\beta_k}
            S_{\beta_{k+1}}\cdots S_{\beta_{m-1}}
            S_{\beta_m}                                \\
             & =
            S_{\beta_1}\cdots S_{\beta_{k-1}}
            \bigl(uS_{\beta_m}u^{-1}\bigr)uS_{\beta_m} \\
             & =
            S_{\beta_1}\cdots S_{\beta_{k-1}}u         \\
             & =
            S_{\beta_1}\cdots S_{\beta_{k-1}}
            S_{\beta_{k+1}}\cdots S_{\beta_{m-1}}.
        \end{aligned}
    \]
    This is a shorter expression for \(w\), contradicting the choice of
    \[
        w=S_{\beta_1}\cdots S_{\beta_m}.
    \]
    Thus
    \[
        w(\beta_m)\in \Delta^-.
    \]

    On the other hand, since \(w(B)=B\) and \(\beta_m\in B\), we have
    \[
        w(\beta_m)\in B\subseteq \Delta^+,
    \]
    a contradiction. Therefore \(w=1\).
\end{proof}

\begin{corollary}
    If \(B\) and \(B'\) are two bases of \(\Delta\), then there is at most one
    element \(w\in W_B\) such that
    \[
        w(B')=B.
    \]
\end{corollary}

\begin{proof}
    Suppose that \(w_1,w_2\in W_B\) satisfy
    \[
        w_1(B')=B
        \qquad\text{and}\qquad
        w_2(B')=B.
    \]
    Then
    \[
        w_2w_1^{-1}(B)=B.
    \]
    By the proposition,
    \[
        w_2w_1^{-1}=1.
    \]
    Hence \(w_1=w_2\).
\end{proof}

\begin{definition}[Normal system]
    Let \(E\) be a Euclidean space. A set of vectors
    \[
        v_1,\dots,v_k\in E
    \]
    is called a \emph{normal system} if the following conditions hold:
    \begin{enumerate}
        \item \(v_1,\dots,v_k\) are linearly independent;

        \item
              \[
                  \|v_i\|=1
                  \qquad\text{for all } i;
              \]

        \item for \(i\neq j\),
              \[
                  (v_i\mid v_j)
                  \in
                  \left\{
                  0,\,
                  -\frac{1}{2},\,
                  -\frac{\sqrt{2}}{2},\,
                  -\frac{\sqrt{3}}{2}
                  \right\}.
              \]
    \end{enumerate}
\end{definition}

\begin{example}
    Let
    \[
        B=\{\beta_1,\dots,\beta_n\}
    \]
    be a base of a root system \(\Delta\). Then
    \[
        \left\{
        \frac{\beta_1}{\|\beta_1\|},
        \dots,
        \frac{\beta_n}{\|\beta_n\|}
        \right\}
    \]
    is a normal system.
\end{example}

\begin{definition}[Coxeter graph associated with a normal system]
    Let
    \[
        \{v_1,\dots,v_k\}
    \]
    be a normal system. Its \emph{Coxeter graph} is defined as follows:
    \begin{itemize}
        \item the vertices are \(v_1,\dots,v_k\);
        \item for \(i\neq j\), put
              \[
                  4(v_i\mid v_j)^2
              \]
              edges between \(v_i\) and \(v_j\).
    \end{itemize}
    Thus two vertices are connected precisely when
    \[
        (v_i\mid v_j)\neq 0.
    \]
\end{definition}

\begin{proposition}[Properties of normal systems and Coxeter graphs]
    Let
    \[
        \{v_1,\dots,v_k\}
    \]
    be a normal system. Then the following hold.
    \begin{enumerate}
        \item Every subset of a normal system is again a normal system.

        \item The number of connected pairs \(v_i-v_j\) is at most \(k-1\).

        \item The Coxeter graph contains no cycles.

        \item Fix \(v_i\). Then the total number of edges adjacent to \(v_i\) is
              strictly less than \(4\). Equivalently, it is at most \(3\).

        \item Suppose \(v_1,\dots,v_k\) form a simple chain in the Coxeter graph:
              \[
                  v_1 - v_2 - \cdots - v_{k-1}-v_k,
              \]
              namely
              \[
                  (v_i\mid v_j)=
                  \begin{cases}
                      -\dfrac{1}{2}, & j=i+1,            \\[4pt]
                      0,             & \text{otherwise}.
                  \end{cases}
              \]
              Let
              \[
                  v=v_1+\cdots+v_k.
              \]
              Replacing \(v_1,\dots,v_k\) by \(v\) gives a new normal system.
    \end{enumerate}
\end{proposition}

\begin{proof}
    \begin{enumerate}
        \item This is immediate from the definition of a normal system.

        \item Let
              \[
                  v=v_1+\cdots+v_k.
              \]
              Since \(v_1,\dots,v_k\) are linearly independent, we have \(v\neq 0\).
              Hence
              \[
                  (v\mid v)>0.
              \]
              On the other hand,
              \[
                  \begin{aligned}
                      (v\mid v)
                       & =
                      \sum_{i=1}^k (v_i\mid v_i)
                      +
                      2\sum_{i<j}(v_i\mid v_j) \\
                       & \leq
                      k-\#\{\text{connected pairs}\},
                  \end{aligned}
              \]
              because if \(v_i\) and \(v_j\) are connected, then
              \[
                  (v_i\mid v_j)\leq -\frac{1}{2}.
              \]
              Therefore
              \[
                  \#\{\text{connected pairs}\}<k.
              \]
              Since this number is an integer, it is at most \(k-1\).

        \item If the Coxeter graph contained a cycle with \(r\) vertices, then the
              corresponding \(r\) vectors would form a normal subsystem whose Coxeter
              graph has at least \(r\) connected pairs. This contradicts \textup{(2)},
              which says that a normal system of \(r\) vectors has at most \(r-1\)
              connected pairs.

        \item Let
              \[
                  v_{j_1},\dots,v_{j_l}
              \]
              be the vertices adjacent to \(v_i\), and let \(k_s\) denote the number of
              edges between \(v_i\) and \(v_{j_s}\). Thus
              \[
                  k_s=4(v_i\mid v_{j_s})^2.
              \]

              By \textup{(3)}, the vertices \(v_{j_1},\dots,v_{j_l}\) must be pairwise
              orthogonal. Otherwise we would obtain a cycle.

              Consider
              \[
                  u
                  =
                  v_i
                  -
                  \sum_{s=1}^l (v_i\mid v_{j_s})v_{j_s}.
              \]
              Since the vectors in a normal system are linearly independent, we have
              \(u\neq 0\). Hence
              \[
                  (u\mid u)>0.
              \]
              Since \(v_{j_1},\dots,v_{j_l}\) are pairwise orthogonal and have norm
              \(1\), we get
              \[
                  \begin{aligned}
                      (u\mid u)
                       & =
                      (v_i\mid v_i)
                      -
                      \sum_{s=1}^l (v_i\mid v_{j_s})^2 \\
                       & =
                      1-\frac{1}{4}\sum_{s=1}^l k_s.
                  \end{aligned}
              \]
              Therefore
              \[
                  \sum_{s=1}^l k_s<4.
              \]
              Hence the total number of edges adjacent to \(v_i\) is strictly less than
              \(4\).

        \item Let
              \[
                  S=\{v_1,\dots,v_k\}\cup T
              \]
              be the original normal system, where \(T\) denotes the remaining vectors.
              Set
              \[
                  v=v_1+\cdots+v_k.
              \]
              We show that
              \[
                  S'=\{v\}\cup T
              \]
              is again a normal system.

              First, \(S'\) is linearly independent. Indeed, if
              \[
                  av+\sum_{w\in T}a_w w=0,
              \]
              then
              \[
                  a(v_1+\cdots+v_k)+\sum_{w\in T}a_w w=0.
              \]
              Since the original system \(S\) is linearly independent, we get
              \[
                  a=0
                  \qquad\text{and}\qquad
                  a_w=0
                  \quad\text{for all }w\in T.
              \]

              Next,
              \[
                  \begin{aligned}
                      (v\mid v)
                       & =
                      \left(v_1+\cdots+v_k\mid v_1+\cdots+v_k\right) \\
                       & =
                      \sum_{i=1}^k (v_i\mid v_i)
                      +
                      2\sum_{i=1}^{k-1}(v_i\mid v_{i+1})             \\
                       & =
                      k+2(k-1)\left(-\frac{1}{2}\right)              \\
                       & =1.
                  \end{aligned}
              \]
              Hence \(\|v\|=1\).

              Finally, let \(w\in T\). By \textup{(3)}, the vertex \(w\) can be connected
              to at most one of \(v_1,\dots,v_k\); otherwise the Coxeter graph would
              contain a cycle. Hence
              \[
                  (w\mid v)
                  =
                  (w\mid v_1)+\cdots+(w\mid v_k)
              \]
              is either \(0\), or equal to one of the inner products \((w\mid v_i)\).
              Therefore
              \[
                  (w\mid v)
                  \in
                  \left\{
                  0,\,
                  -\frac{1}{2},\,
                  -\frac{\sqrt{2}}{2},\,
                  -\frac{\sqrt{3}}{2}
                  \right\}.
              \]
              All other inner products are unchanged. Thus \(S'\) is a normal system.
    \end{enumerate}
\end{proof}

